% ৩.৪ বাইনারি যোগ ও বিয়োগ
\section{বাইনারি যোগ ও বিয়োগ}

\subsection{Binary Addition}
\begin{itemize}
\item 0+0=0, 0+1=1, 1+1=10 (carry)
\item Carry ধারণা ও multi-bit যোগ
\end{itemize}

\subsection{Binary Subtraction}
\begin{itemize}
\item 0-0=0, 1-0=1, 1-1=0, 0-1= borrow
\item Borrow ধারণা
\end{itemize}

\subsection{প্রশিক্ষণ সমস্যা}
\begin{itemize}
\item পরীক্ষার ধাঁচের উদাহরণ
\end{itemize}

\begin{learningbox}
এই টপিক শেষে শিক্ষার্থী পারবে:
\begin{itemize}
\item binary addition ও subtraction-এর basic rule লিখতে।
\item carry ও borrow ব্যবহার করে multi-bit অংক করতে।
\item binary arithmetic-এ common mistake এড়াতে।
\end{itemize}
\end{learningbox}

\begin{center}
\begin{tabular}{|p{0.25\textwidth}|p{0.25\textwidth}|p{0.30\textwidth}|}
\hline
\textbf{Addition} & \textbf{ফল} & \textbf{নোট} \\
\hline
$0+0$ & 0 & carry 0 \\
\hline
$0+1$ বা $1+0$ & 1 & carry 0 \\
\hline
$1+1$ & 10 & sum 0, carry 1 \\
\hline
$1+1+1$ & 11 & sum 1, carry 1 \\
\hline
\end{tabular}
\end{center}

\begin{examplebox}{Binary যোগ}
\[
\begin{array}{r}
  1011\\
+ 0110\\ \hline
 10001
\end{array}
\]
এখানে ডান দিক থেকে bit যোগ করে carry পরের column-এ নেওয়া হয়েছে।
\end{examplebox}

\begin{examplebox}{Binary বিয়োগ}
\[
\begin{array}{r}
  10110\\
- 01011\\ \hline
  01011
\end{array}
\]
যেখানে $0-1$ এসেছে, সেখানে বাম পাশের bit থেকে borrow নেওয়া হয়েছে।
\end{examplebox}

\begin{mistakebox}{সাধারণ ভুল}
Binary addition-এ $1+1=2$ লেখা যাবে না; binary-তে 2 প্রকাশ হয় $10$ হিসেবে। তাই ফল 0 এবং carry 1।
\end{mistakebox}
